% !PRINTED 250
\weberpage{250}
%<<CONT
two imaginary roots, and the points of $[2, 4]$ equations with two equal roots and two real roots distinct from them.

On the double parabola two equal roots occur twice, and indeed on the part where $[0]$ borders on $[4]$ they are real, on the isolated part imaginary.

% Source figure: Weber_Algebra_Bd1_IA.pdf, PDF page 276, printed p.250, Fig.5; exact caption-inclusive crop.
% Former placeholder description retained: perspective line drawing of the discriminant surface, showing two cuspidal-edge branches meeting at a shaded curved vault-like shell; dotted lines indicate hidden edges and the coordinate framework of the model.
\weberfigureimage{b1_i0275_fig5_source_20260831.png}{0.56}

The points of the cuspidal edges represent equations with three equal roots, and the origin of coordinates the equation with four equal roots.

It was \gesperrt{Kronecker} who first drew attention to the discriminant surface and its significance for the discussion of the biquadratic equation. A model of the surface has been constructed by \gesperrt{Kerschensteiner}.\footnote{\gesperrt{Kronecker}, Monatsbericht der Berliner Akademie of 14 February 1878. A description of the model is to be found in the catalogue of mathematical models edited by Dyck, Munich 1892, from which Fig. 5 above is taken.}

The analytic expression for the conditions of reality of the roots can be brought into a form in which only the covariants of the biquadratic form occur.\footnote{\gesperrt{Clebsch}, ``Theorie der binären Formen'' \S.\,47. \gesperrt{Faà di Bruno}, ``Theorie der binären Formen'', German translation by Walter (Leipzig 1881), \S.\,20, where an essential addition by \gesperrt{Nöther} is to be found.}

The starting point for this is furnished by the expressions for the functions $\psi_1$, $\psi_2$, $\psi_3$ given in \S.\,66, (6). If we, in
%CONT>>
